\section{Related Work}
\label{sec:relatedWork}

Research on business rules in software engineering has addressed their definition, characterization, and extraction from software artifacts. These tasks differ in their objectives: characterization describes the nature and structure of business rules, extraction aims to recover rules from software artifacts, whereas localization starts from a known rule and seeks to identify its implementation in source code. Comparatively less attention has been given to locating the implementation of a known business rule in source code.

In this section, we first review existing definitions and characterizations of business rules. We then examine how business vocabulary is represented in source code through business classes and variables, and how these elements relate to the identification of business-related code. Next, we review approaches for extracting business rules from source code, database metadata, and, more recently, using large language models (LLMs). We then discuss the limitations of existing evaluation practices before examining research on business-rule localization, with a focus on documentation-to-source-code traceability. Finally, we highlight the limitations of existing approaches in locating the method implementing a business rule from its textual description.

\subsection{Business Rule Definitions and Characterization}

To understand business rules, it is necessary to examine how they are defined and characterized.

Business rules are operational rules that guide organizational activities and define constraints, decisions, or conditions governing system behavior~\cite{Huan96a}. They may originate from legal requirements, organizational policies, or industry practices and are typically implemented within application logic~\cite{Wang08b}. In software systems, they can be scattered across classes and methods, making them difficult to identify and maintain~\cite{Cose13a}.

The characterization of business rules has been studied since the 1990s. Sneed and Erdos~\cite{Snee96a} proposed a canonical structure composed of results, arguments, assignments, and conditions. Wang \etal~\cite{Wang08b,Wang04b,Wang04a} later defined business rules as conditional operations associated with specific data results. Chiang~\cite{Chia06a} characterized a primitive business rule as a function that takes inputs and produces an output, while Gang~\cite{Gang09a} considered business rules as functions, constraints, or transformations from inputs to outputs.

These definitions provide useful abstractions for identifying business logic in source code. However, they are generally broad, and not every function or condition necessarily represents business logic. In large systems, business rules may therefore be difficult to distinguish from technical implementation code.

\subsection{Business Vocabulary and Code Elements}
\label{sec:businessVar}

Understanding how business vocabulary is represented in source code helps identify relevant code elements.

A business vocabulary captures the concepts and relationships of a business domain. Prakash \etal~\cite{Prak21a} define a business vocabulary as a set of interconnected concepts consisting of \textit{business terms}, which represent business entities and concepts, and verbs, which express relationships between them. Examples of business terms include \textit{Currency}, \textit{Customer}, \textit{Operation} and \textit{Operational Cost}~\cite{Ayac26a}.

\textit{Business terms} can be represented in source code through \textit{business variables}, which provide concrete representations of business vocabulary~\cite{Wang08b}, and \textit{business classes}, which define the structure and behavior of the corresponding objects in an object-oriented context~\cite{Ayac26a}. For example, the business term \textit{operation} can be represented by the class \texttt{Operation}, while variables such as \texttt{op1} and \texttt{op2} represent instances of this class. Identifying these source-code elements can help understand how business concepts, and ultimately business rules, are represented and implemented in a software system.

Several studies have investigated the identification of business variables by exploiting their relationships with the data manipulated by the system. Wang et al.~\cite{Wang08b} consider this identification as an initial step in their approach and focus on variables involved in system inputs and outputs. They argue that when the database is regarded as a source of input and output data, its fields are likely to contain values corresponding to business variables. Similarly, Sneed and Verhoef~\cite{Snee21a} trace output variables back to their origins in the program and show that some of these values can be linked directly to database tables. These studies highlight the relevance of database information for identifying business-related elements in source code. However, manually identifying such elements can be difficult in large industrial systems, where numerous variables must be examined, often with the support of domain experts.

In object-oriented systems, business classes offer an alternative starting point. Ayachi et al.~\cite{Ayac26a} explored this approach by focusing on persistent business classes, which are mapped to database tables through the Object-Relational Mapping (ORM) layer. This relationship provides a structural link between business data and its representation in source code, which can be exploited to identify code elements relevant to business rules.

\subsection{Extraction of Business Rules}

Business rules can be extracted from various software artifacts using different techniques.

\subsubsection{Extraction of Business Rules from Source Code}\noindent

Several approaches have investigated the extraction of business rules directly from source code. Sneed and Erdos~\cite{Snee96a} proposed tracing backwards from an identified result to reconstruct the corresponding assignments, arguments, and conditions. Wang \etal~\cite{Wang08b,Wang04b,Wang04a} focused on identifying conditional operations associated with business variables. Chiang~\cite{Chia06a} and Gang~\cite{Gang09a} used program slicing to isolate portions of code relevant to business rules.

Although these approaches differ in their analysis techniques, they generally require an initial identification of relevant business variables or output points. This manual step is difficult in large systems containing thousands of variables. Moreover, a substantial part of the identified operations may correspond to technical rather than business logic~\cite{Snee20b}.

\subsubsection{Extraction of Business Rules from Database Metadata}\noindent

Business rules have also been extracted from database artifacts. Chaparro \etal~\cite{Chap12a} investigated the extraction of structural business rules from databases by exploiting their business vocabulary, including table and attribute names. Their work highlights that database structures can provide useful information about the business concepts represented in a software system.

\subsubsection{Extraction of Business Rules Using LLMs}\noindent

More recently, large language models have been used to extract business rules from source code. Miskell \etal~\cite{Misk23a}, for example, proposed an approach that uses LLMs to analyze Java classes and methods and generate business rules expressed as natural-language sentences. While such approaches demonstrate the potential of LLMs for recovering business knowledge from code, they start from source code and attempt to recover the corresponding rules. Our work addresses the inverse problem: given a textual description of a known business rule, locating its implementing method.

\subsection{Limitations in Evaluating the Approaches}
\label{sec:limitations}

Despite the progress made in business-rule extraction, evaluating these approaches remains challenging. Several studies~\cite{Chia06a,Snee96a} do not report an empirical evaluation, while others~\cite{Wang04b,Wang04a,Snee20b,Chap12a} use basic indicators, such as the number of rules or business-related elements identified. However, the actual number of business rules present in a software system is rarely known, making it difficult to assess how completely an approach covers the existing rules. Expert-based validation~\cite{Gang09a} can provide valuable evidence, but manually examining a large set of extracted rules requires substantial effort. More recent approaches using LLMs~\cite{Misk23a} introduce another consideration: reproducibility. Their outputs may vary across executions, even when the same input is provided, which can make the results of an evaluation less consistent.

\subsection{Locating Business Rules}

While characterization and extraction have been studied, the explicit localization of business rules in source code remains underexplored and challenging in complex systems. One related area is documentation-to-source-code traceability, which aims to establish links between software documentation and the corresponding code elements.

Traceability between software documentation and source code has been widely studied to help developers find the code related to requirements and other documented information. Existing approaches use different techniques, such as information retrieval, machine learning, and source-code analysis.

Antoniol et al.~\cite{Anto02b} proposed an information retrieval approach to link source code with free-text documentation. Their method uses code identifiers, such as class and method names, to search for relevant documents and rank them according to their textual similarity. Their results showed that this approach can identify relevant links while reducing the number of documents that developers need to examine during the search process.

Other studies have tried to improve these links by using additional information. Cao et al.~\cite{Cao18a} used machine learning to distinguish code elements that are important to a document from those mentioned only as context. Their approach uses textual and structural information to identify more relevant links. Ali et al.~\cite{Ali11a} proposed Coparvo, which separates source code into different parts, such as class names, method names, variable names, and comments. These parts help assess the links initially identified through information retrieval and reduce false positives.

Some approaches combine several techniques to improve the results. Eaddy et al.~\cite{Eaddy08a} proposed CERBERUS, which combines information retrieval with dynamic analysis and program dependency analysis to identify code related to software requirements. More recently, Keim et al.~\cite{Keim24a} proposed TransArC to link architectural documentation to source code through an intermediate architecture model. Their approach uses information such as naming similarities, package structures, and inter-component relationships.

Although these approaches help developers link documentation to source code, they address a different problem and do not directly fit our needs. Most of them aim to link general requirements or documents to code elements. In our case, the task is more specific: the input consists of the description of a single business rule, which may be expressed in just one sentence or a few lines of text. This provides only a limited amount of textual information to guide the search, compared with approaches that rely on larger documentation sources containing more extensive descriptions and contextual information. Moreover, we aim to identify the specific implementation of the rule, rather than a broader set of related code elements. Our approach therefore operates at a finer level of granularity and must identify the implementation using limited textual information.